%%%%%%%%%%%%%%%%%%%%%%%%%%%%%%%%%%%%%%%%%%%%%%%%%%%%%%%%%%%%%%%%%%%%%%%%%%%%%%%%
%   Copyright 2009, Oracle and/or its affiliates.
%   All rights reserved.
%
%
%   Use is subject to license terms.
%
%   This distribution may include materials developed by third parties.
%
%%%%%%%%%%%%%%%%%%%%%%%%%%%%%%%%%%%%%%%%%%%%%%%%%%%%%%%%%%%%%%%%%%%%%%%%%%%%%%%%

\newpage
\section{\wherecore}
\applabel{where-calculus}

\note{Where clauses are not yet supported.}

In this section, we define a Fortress core calculus with \wc.
We call this calculus \emph{\wherecore}.
\wherecore\ is an extension of \basiccore\ with \wc.
\revision{revival-calculi}{This calculus predates instantiation exclusion
(\secref{trait-types}, \secref{trait-decls}) and the rule that a
\KWD{where}-clause variable may not appear as a static argument in an
\KWD{extends} clause (\secref{where-clauses}).
It lets a \KWD{where}-clause variable occur in a supertype of a trait or
object definition, and its subtyping rule \sBothRule\ then makes the
definition a subtype of one instantiation of that supertype for each
witness of the variable: the shape of the declarations marked
``Not allowed'' in \secref{where-clauses}.
The calculus, its rules and its soundness claim are as the team wrote and
proved them, and are not changed.}

\input{\macros/where-calculus-macros}
\input{\calculi/where/syntax}
\input{\calculi/where/dynamic}
\input{\calculi/where/static}
